\documentclass[11pt,a4paper]{article}
\usepackage[T1]{fontenc}
\usepackage[utf8]{inputenc}
\usepackage{amsmath,amssymb,xcolor}
\usepackage[margin=25mm]{geometry}
\usepackage[colorlinks=true,urlcolor=blue]{hyperref}
\definecolor{nexorablue}{HTML}{2459B5}
\setlength{\parindent}{0pt}
\setlength{\parskip}{7pt}
\setlength{\emergencystretch}{2em}
\renewcommand{\familydefault}{\sfdefault}
\begin{document}

{\small\bfseries\color{nexorablue} NEXORA CPGE}\par
{\small MPSI / NEX-CIR-02 / Énoncé}\par
{\LARGE\bfseries Deux sources opposées}\par
\vspace{0.25em}\hrule\vspace{0.6em}

Dans une maille série, le parcours choisi traverse une source de $18\,\mathrm V$ du pôle négatif au positif, une résistance de $4\,\mathrm\Omega$, puis une source de $6\,\mathrm V$ du pôle positif au négatif et une résistance de $2\,\mathrm\Omega$. Orienter le courant suivant ce parcours. Donner son signe, les quatre puissances reçues et le rôle des deux sources.


\end{document}
